%!TEX ROOT=sample-thesis.tex
%!TEX PROGRAM=xelatex

\chapter{Problem Characterization Interview Sample}

\begin{table}[hbt!]
\centering
\begin{tabular}{|p{3cm}|p{9.5cm}|} % max 12.5cm
	\hline
	Questions & Voice remarks (\textcolor{blue}{Researcher = Blue}, Experts = Dark) \\
  	\hline
  	\multicolumn{2}{|c|}{Problem Characterization} \\
  	\hline
	Q1 & Okay, mine is basically task-based, it was not content-based. all the modules are pretty much task-oriented, so there are tasks to be completed. if we look at the old version, it looks similar to the forum when the learner completed the tasks. the form is like there is a video, and there are discussions attached to it. because mine is self-paced, people respond to the tasks, it is not meant for them to discuss in the forum, it is meant for them to go to the tasks, it just that it is more task-based. however, if somebody wants to read other people's comments, they want to comment, it is there. at a few levels, there are peers commenting on their tasks, that is it. but the whole design is not a forum. the whole design is for each of them to complete a task, so the task sometimes requires them to answer in sentences. but many of my tasks were also in a form of visual. they need to create and submit a poster, a mind map. so, the interaction happens there is like peers upvoting or 'like' their work. \newline
	\textcolor{blue}{That one is one of the activities that we intend to identify such as 'like', 'upvote'.} \\
  \hline
\end{tabular}
\end{table}
